\section{Introduction}

\subsection{Purpose}

This document provides a detailed description of the software design for the medical clinic management system, in accordance with the Software Design Document standard.
It serves as a primary technical reference that explains how the requirements specified in the Software Requirements Specification are transformed into a clear architectural design and implementable software modules.
This document is intended for developers, software engineers, and any technical stakeholders involved in building, maintaining, or evolving the system. It describes the system architecture, data design, module design, and the interaction between system components.


\subsection{Scope}

This design covers a comprehensive web-based system dedicated to managing a medical clinic. The system enables the management of patients, doctors, reception staff, and administrators, while securely maintaining electronic medical records in an organized manner.
The system provides an integrated set of operational functionalities covering user management, patient management, appointment scheduling, medical records, and financial billing, as follows:


\begin{itemize}
	\item \textbf{Priority Categorization:} Every functional requirement is assigned a priority level to guide the development and testing phases:
	\begin{itemize}
		\item \textbf{[High]}: Critical features required for the core operation of the clinic.
		\item \textbf{[Medium]}: Important features that enhance efficiency but are not vital for basic operations.
		\item \textbf{[Low]}: Optional or ``nice-to-have'' features that may be deferred to future updates.
	\end{itemize}
	
	\item \textbf{Typographical Styles:}
	\begin{itemize}
		\item \textbf{Boldface:} Used for Graphical User Interface (GUI) elements, such as buttons, menu items, and window titles (e.g., ``Click the \textbf{Save} button'').
		\item \textit{Italics:} Used for emphasis, key technical terms at their first mention, or placeholders for specific values.
		\item \texttt{Monospace Font:} Used for database table names, field names, or code snippets (e.g., \texttt{Patient\_ID}).
	\end{itemize}
	
	\item \textbf{Identification Schemes:} Standard prefixes are used for cross-referencing between sections:
	\begin{itemize}
		\item \texttt{REQ-x}: Represents a Functional Requirement (e.g., \texttt{REQ-1}).
		\item \texttt{UC-x}: Represents a Use Case (e.g., \texttt{UC-05}).
		\item \texttt{TBD-x}: Represents a placeholder for items ``To Be Determined'' (e.g., \texttt{TBD-01}).
	\end{itemize}
\end{itemize}

\subsection{Intended Audience and Reading Suggestions}

This document is intended for the following audience:

\begin{itemize}
	\item \textbf{Stakeholders \& Clinic Administration:} To understand how the system addresses current operational challenges and ensures alignment with the clinic’s strategic goals.
	\item \textbf{System Analysts:} To utilize the Business Rules and Functional Requirements as a foundation for designing UML diagrams and system logic.
	\item \textbf{Developers:} To translate the defined Functional Requirements and Database Schema into efficient and accurate source code.
	\item \textbf{Quality Assurance (QA) Team:} To use the described scenarios and features as a primary reference for building test plans and ensuring the system is bug-free before delivery.
\end{itemize}

\subsection{Product Scope}

\subsubsection{General Description}

The \textit{Clinic Management System} is an integrated administrative and medical platform designed to digitize and automate manual clinic operations. The system connects four key stakeholders (\textit{Patient, Receptionist, Doctor, and Administration}) into a single data loop, ensuring a seamless flow of information from initial appointment booking to final billing and medical record archiving.

\subsubsection{In-Scope (Included Functional Requirements)}

Based on the defined database schema, the project will deliver the following software solutions:

\begin{itemize}
	\item \textbf{Smart Booking System:} Automating appointments and performing real-time conflict checks against doctor schedules.
	\item \textbf{Digital Medical Records (EMR):} Converting physical patient files into electronic records, including diagnoses and prescriptions.
	\item \textbf{Access Control \& Auditing:} Implementing Role-Based Access Control (RBAC) and maintaining an Active Log to track all system transactions and prevent unauthorized tampering.
	\item \textbf{Financial Engine:} Automated calculation of medical service costs and invoice generation (\textit{Bills}).
	\item \textbf{Specialty Management:} Managing medical specialties to streamline patient referrals to the appropriate physician.
\end{itemize}

\subsubsection{Out-of-Scope (Excluded from Current Phase)}

To ensure timely delivery and focus on core objectives, the following features are excluded from the current version:

\begin{itemize}
	\item Internal Pharmacy System.
	\item External Lab Integration.
	\item Mobile Application for patients.
	\item Human Resources \& Payroll management.
\end{itemize}

\subsection{References}

Reference materials and standards:

\begin{itemize}
	\item IEEE Std 830-1998: IEEE Recommended Practice for Software Requirements Specifications.
	\item Internal clinic documentation: patient record templates, booking logs, and financial receipt samples.
	\item User interview notes gathered from receptionists and doctors regarding current workflow and needs.
\end{itemize}